\documentclass[a4paper,11pt]{article}
\usepackage[utf8]{inputenc}
\usepackage[T1]{fontenc}
\usepackage[english]{babel}
\usepackage[left=2.5cm,right=2.5cm,top=2cm,bottom=2cm]{geometry}
\usepackage{hyperref}
\usepackage{xcolor}
\usepackage{microtype}
\usepackage{lmodern}

\definecolor{primary}{RGB}{0, 45, 114}
\hypersetup{colorlinks=true, urlcolor=primary}

\begin{document}
\noindent
\textbf{Loïc Aron MBASSI EWOLO} \\
ENSEA -- Double Diplôme Ingénieur \\
\href{mailto:loicaron.mbassiewolo@ensea.fr}{loicaron.mbassiewolo@ensea.fr} \quad (+33) 7 59 52 33 98 \\
\href{https://github.com/Nameless0l}{github.com/Nameless0l} \quad \href{https://mbassiloic.tech}{mbassiloic.tech}

\vspace{1em}
\noindent Cergy, le 22 mars 2026

\vspace{1em}
\noindent
\textbf{Objet : Candidature -- Agents IA Cybersécurité} \\
\textbf{Référence : Offre d'alternance -- Thales}

\vspace{1.5em}
\noindent Madame, Monsieur,

\vspace{0.8em}
Les systèmes de défense moderne exigent une IA capable de raisonner, d'apprendre et de s'adapter en temps réel. Vous recherchez un ingénieur capable de concevoir des agents IA orchestrés pour la cybersécurité. Mon expérience en systèmes multi-agents (Google ADK), en pipeline NLP de recherche (INRIA Saclay) et en architecture défensive (protocoles Signal, analyse de confiance) me permet de designer et implémenter rapidement des agents autonomes pour vos défenses critiques.

\vspace{0.8em}
Durant mon stage à INRIA Saclay (équipe TRIBE, avril-août 2026), je mets en œuvre un pipeline NLP avancé pour analyser et évaluer automatiquement la confiance des applications mobiles. Ce travail m'expérimente aux défis de l'automatisation d'analyses complexes : orchestration d'étapes (détection, extraction, scoring), fusion de signaux hétérogènes, génération de rapports interpétables. J'ai parallèlement conçu un système multi-agents agricole utilisant Google ADK, où cinq agents collaboratifs (météo, marché, sols, cultures, recommandation) orchestrés via Gemini produisent des décisions cohérentes. Cette architecture m'a enseigné comment structurer le raisonnement distribué et gérer les conflits entre agents.

\vspace{0.8em}
Mon expertise en architecture sécurisée provient du projet Snappy : implémentation du protocole Signal (Perfect Forward Secrecy) et cryptage AES-256-GCM combiné à une évaluation RGPD complète du système. Cette expérience m'a montré comment encoder les politiques de sécurité en logique métier : qui a accès à quoi, quand, et comment s'en assurer. Le déploiement sur CoHoMa (robotique militaire autonome) m'a formé aux contraintes du temps réel embarqué, où les agents doivent décider et agir rapidement malgré les ressources limitées (GPU Nvidia Jetson).

\vspace{0.8em}
Je suis disponible dès septembre 2026 pour une alternance de 12 mois. Convaincu que les défenses cybermodernes reposent sur des agents IA raisonnants et adaptatifs, je suis prêt à concrétiser votre vision des systèmes autonomes sécurisés.

\vspace{1.5em}
\noindent Veuillez agréer, Madame, Monsieur, l'expression de mes salutations distinguées.

\vspace{2em}
\noindent \textbf{Loïc Aron MBASSI EWOLO}

\end{document}
